\documentclass[11pt,a4paper]{article}
\usepackage[margin=24mm,headheight=15pt]{geometry}
\usepackage{fontspec}
\setmainfont{DejaVu Serif}
\setsansfont{DejaVu Sans}
\setmonofont{DejaVu Sans Mono}[Scale=0.82]
\usepackage{babel}
\babelprovide[main,import]{russian}
\setlength{\emergencystretch}{2em}
\usepackage{amsmath,amssymb,amsthm,mathtools}
\usepackage{booktabs,tabularx,array,enumitem,xcolor,fancyhdr,hyperref,microtype}
\definecolor{ink}{HTML}{183C4A}
\definecolor{accent}{HTML}{286D70}
\hypersetup{colorlinks=true,linkcolor=accent,urlcolor=accent,pdftitle={D0: одно проективное масштабирование и жёсткость источника Эйнштейна},pdfauthor={Исследовательская сессия D0, 18 ноября 2026}}
\pagestyle{fancy}
\fancyhf{}
\fancyhead[L]{\small\sffamily D0 / исследовательская записка}
\fancyhead[R]{\small\sffamily 18.11.2026}
\fancyfoot[C]{\small\thepage}
\setlength{\parindent}{0pt}
\setlength{\parskip}{6pt}
\setlist{itemsep=4pt,topsep=5pt,leftmargin=*}
\renewcommand{\arraystretch}{1.15}
\newcommand{\E}{\mathcal E}
\newcommand{\dd}{\mathrm d}
\newcommand{\R}{\mathbb R}
\DeclareMathOperator{\Ric}{Ric}
\DeclareMathOperator{\Tr}{Tr}
\DeclareMathOperator{\Span}{span}
\newtheorem{theorem}{Теорема}
\newtheorem{lemma}{Лемма}
\newtheorem{corollary}{Следствие}
\begin{document}
\sloppy
{\sffamily\color{accent}\large D0 $\longrightarrow$ ОТО}
\vspace{5mm}

{\sffamily\color{ink}\LARGE Одно проективное масштабирование\\[2mm]
и жёсткость источника Эйнштейна}

\vspace{4mm}
{\large Полное доказательство, условия применения и проверка физического перехода C}
\vspace{4mm}

\textbf{Основной результат.} Для локального естественного метрического источника точное равенство
$\E[\lambda^2 g]=\E[g]$ можно заменить существенно менее конкретным условием
\[
\boxed{\E[\lambda^2g](x)=r[g](x)\E[g](x),\qquad r[g](x)\ne0.}
\]
Если источник имеет ненулевой квадратичный отклик Эйнштейна около плоской метрики и удовлетворяет гравитационному тождеству сохранения до наложения уравнений движения, то при перечисленных ниже условиях
\[
\boxed{\E[g]=aG[g],\qquad
G_{\mu\nu}=\Ric_{\mu\nu}-\tfrac12 Rg_{\mu\nu},\quad a\ne0.}
\]
Достаточен \emph{один} фиксированный масштаб $\lambda>1$. Формула, знак и гладкость множителя $r$ не требуются. Ограничение на дифференциальный порядок источника не вводится.

\textbf{Научный статус.} Это условная аналитическая теорема, полученная в данной исследовательской сессии и прошедшая две внутренние независимые проверки. Приведены точные символьные контроли. Формализация в Lean и внешнее научное рецензирование не выполнены; литературный приоритет не установлен.

\textbf{Физический статус.} Теорема не доказывает, что реальная полная цена D0 уже порождает такой источник. Обязательство C остаётся открытым; \texttt{GR\_closed=false}. Пользователь разрешил возобновить работу над C в этой сессии. Эмпирические блоки не исследовались.

\section{Исходная точка и точное изменение}
Исследование опирается на проверенную ветку PR~334 репозитория \texttt{gvakhrushev/d0\_15}, commit
\begin{center}
\texttt{dc411daf01f697ec85dd22f28bfe124171796645}.
\end{center}
В основном научном документе раздел \S7fu сохраняет математическую композицию полного источника и условное восстановление сектора Einstein--two-CAR, но отзывает утверждение о доказанном физическом переходе C~\cite{primary}. Совпадение операторной статистической суммы с некоторой графовой конструкцией ещё не является переносом действительной цены \emph{той же полной истории}.

Ранее имелся критерий единственности с точным однократным масштабным законом. Его приложение G12c уже учитывало исчезающие ошибки конечного приближения~\cite{golden}. Поэтому утверждение «достаточно приближённого масштабирования в пределе» не было бы новым результатом этой работы.

Здесь меняется сам достаточный критерий: требуется сохранение \emph{одномерного направления тензора}, без заранее вычисленной нормировки. Это не обычное совпадение решений на одной семье волн: условие проверяется на всей объявленной области метрик, включая вариации, не удовлетворяющие уравнениям движения.

Классическая теорема Navarro--Sancho характеризует естественный сохраняющийся тензор Эйнштейна через вес ноль при всех изменениях масштаба~\cite{navarro}. Настоящая записка даёт отдельное доказательство проективного критерия относительно уже имевшегося D0-критерия. Совпадения формулировки с ранее опубликованными результатами не исключены.

\section{Точная формулировка}
Работаем в четырёх измерениях с гладкими лоренцевыми метриками. Источник $\E$ сопоставляет метрике гладкий симметричный ковариантный двухтензор. Все предпосылки относятся к \emph{одному и тому же} источнику.

\begin{enumerate}
\item \textbf{Локальность по ростку и естественность.} Значение источника около точки определяется метрикой в сколь угодно малой окрестности этой точки. Операция коммутирует с ограничением области и с локальными диффеоморфизмами:
\[
\E[\phi^*g]=\phi^*\E[g].
\]
Сохранившиеся независимые часы, записи, внешняя метрика или глобальная история не являются скрытыми аргументами $\E$.

\item \textbf{Открытая область вне уравнений движения.} Допускаются необходимые дискретные масштабирования $\lambda^{2n}g$, их ограниченные обратные образы, нормальные сжатия
\[
g_t(x)=g(tx),\qquad 0<t\le1,
\]
а также достаточно малые гладкие локальные возмущения и склейки, используемые в доказательстве. Допуск не обусловлен заранее равенством $G=0$. Ограниченный физический chart, не содержащий этих операций, автоматически теоремой не покрывается.

\item \textbf{Регулярность.} Источник непрерывен из компактно-открытой $C^\infty$-топологии метрик в непрерывные тензорные поля, в том числе на искривлённых метриках. На постоянных лоренцевых метриках он дифференцируем по Фреше; для вывода плоской леммы достаточно дифференцируемости в $C^1$.

\item \textbf{Плоский вакуум.} $\E[\eta]=0$ для постоянных лоренцевых метрик $\eta$.

\item \textbf{Гравитационное тождество Уорда.}
\[
\nabla_g^\mu \E[g]_{\mu\nu}=0
\]
тождественно на всей допустимой области, до наложения уравнений движения.

\item \textbf{Ненулевой квадратичный плоский отклик.} В нормальных координатах около $p=0$ положим
\[
h_2=\tfrac12 g_{ij,kl}(0)x^kx^l\,\dd x^i\dd x^j.
\]
Требуется один универсальный коэффициент $a\ne0$, для которого
\[
D\E_\eta[h_2](0)=aG[g](0).
\]
Естественность и линеаризованное тождество Уорда вынуждают именно эйнштейновскую форму этого отклика. Ненулевость $a$ требует отдельного свидетеля; одной ненулевости полного плоского дифференциала недостаточно.

\item \textbf{Один проективный масштаб.} Для одного фиксированного $\lambda>1$, для всякой допустимой метрики и точки существует конечное ненулевое число $r[g](x)$:
\[
\E[\lambda^2g](x)=r[g](x)\E[g](x).
\]
Знак, формула, локальность и гладкость $r$ не предполагаются.
\end{enumerate}

\begin{theorem}[Проективная жёсткость]
При этих предпосылках $\E[g]=aG[g]$ на всей объявленной области, включая точки с вырожденным или нулевым тензором Эйнштейна. В частности, $r[g](x)=1$ везде, где $G[g](x)\ne0$.
\end{theorem}

\textbf{Материя.} Для полной системы metric--matter тождество Нётер обычно содержит материальные Euler-строки:
\[
2\nabla_\mu E_g^{\mu\nu}=\sum_A E_{\psi_A}\nabla^\nu\psi_A
\]
с соответствующей конвенцией знаков. Теорема относится к действительно выделенной гравитационной части либо к обоснованному ограничению с нулевой материей. Материальные уравнения нельзя молча считать выполненными в рассуждении вне решений.

\section{Полное доказательство}

\subsection{Нормальное сжатие фиксирует направление}
Выберем нормальные координаты $p=0$ и положим $t_n=\lambda^{-n}$,
$d_t(x)=tx$. Тогда
\[
g_n=t_n^{-2}d_{t_n}^*g=g(t_nx),\qquad
g_n=\eta+t_n^2h_2+o_{C^\infty}(t_n^2).
\]
Естественность и последовательное применение одного проективного закона дают
\[
\E[g_n](0)=t_n^2R_n[g](0)\E[g](0),\qquad
R_n[g](0)=\prod_{k=0}^{n-1}r[\lambda^{2k}g](0).
\]
Поэтому все векторы $t_n^{-2}\E[g_n](0)$ лежат в одном подпространстве
$\Span\{\E[g](0)\}$. Плоская дифференцируемость даёт
\[
t_n^{-2}\E[g_n](0)\longrightarrow aG[g](0).
\]
Одномерное подпространство конечномерного пространства тензоров замкнуто. Отсюда
\[
\E[g](p)\wedge G[g](p)=0.
\]
Здесь клин обозначает линейную зависимость двух элементов $\operatorname{Sym}^2T_p^*M$, а не внешнее произведение пространственно-временных форм.

Если $G[g](p)\ne0$, то $\E[g](p)\ne0$: иначе вся сходящаяся последовательность была бы нулевой. Следовательно, на открытой области $G\ne0$ существует гладкая ненулевая функция $f_g$:
\[
\E[g]=f_gG[g].
\]
Она определяется отношением любой ненулевой компоненты; определения согласуются на пересечениях. Ограниченность $R_n$ и существование его предела нигде не потребовались.

\subsection{Уорд делает коэффициент постоянным в пространстве}
На связной области, где $G$ обратим как билинейная форма,
\[
0=\nabla^\mu(f_gG_{\mu\nu})
  =G_{\mu\nu}\nabla^\mu f_g.
\]
Здесь использовано тождество Бианки для $G$. Обратимость даёт $\dd f_g=0$.

Это ещё не доказывает независимость коэффициента от самой метрики. Разные метрики могли бы иметь разные пространственно постоянные коэффициенты. Следующий шаг устраняет именно эту возможность.

\subsection{Локальная склейка устраняет зависимость от метрики}
Пусть $\det G[g](p)\ne0$. Сузим нормальную координатную область до малого связного шара $U$ так, чтобы $G[g](tx)$ оставался обратимым для $x\in U$ и $0<t\le1$. Для известного тензора Эйнштейна точно
\[
G[g_t](x)=t^2G[g](tx).
\]
Это свойство \emph{сравниваемого} тензора $G$, а не предполагаемая непрерывная масштабная симметрия $\E$.

Зафиксируем $t>0$, точку $q\ne p$ и гладкую функцию $\chi$, равную единице около $p$ и нулю около $q$. Для достаточно малого $\delta$ построим
\[
\widetilde g=g_t+\chi\bigl(g_{t+\delta}-g_t\bigr).
\]
На компактной связной подобласти, содержащей обе точки и носитель $\chi$, это малое $C^2$-возмущение. Метрика остаётся лоренцевой и допустимой, а $G[\widetilde g]$ остаётся обратимым. Поэтому $f_{\widetilde g}$ пространственно постоянно.

Локальность в точках $p$ и $q$ теперь даёт
\begin{align*}
f_{g_{t+\delta}}(p)
&=f_{\widetilde g}(p)
 =f_{\widetilde g}(q)\\
&=f_{g_t}(q)=f_{g_t}(p).
\end{align*}
Таким образом, $t\mapsto f_{g_t}(p)$ локально постоянно на $(0,1]$, следовательно, постоянно на всём этом интервале. Из плоского предела и
$G[g_t](p)=t^2G[g](p)$ получаем $f_{g_t}(p)\to a$, то есть $f_g(p)=a$.

Континуум значений $t$ использован только для допустимых локальных сравнений. Масштабной симметрией $\E$ считался исключительно один дискретный шаг.

\subsection{Вырожденные и нулевые тензоры}
При фиксированных значении и первом джете метрики произвольные достаточно малые изменения симметричного тензора Риччи реализуются квадратичными изменениями метрического джета. Отображение
\[
\Ric\longmapsto \Ric-\tfrac12(\Tr_g\Ric)g
\]
обратимо в размерности четыре. Поэтому ростки с обратимым $G(p)$ плотны.

Берём компактно поддержанные квадратичные возмущения $g_j\to g$ в $C^\infty$ с $\det G[g_j](p)\ne0$. Уже доказано $\E[g_j](p)=aG[g_j](p)$. Непрерывность обоих источников даёт требуемое равенство в пределе. Деления на нулевой тензор нет. Теорема доказана. \hfill$\square$

\section{Более слабый критерий для вакуумных уравнений}
Если ненулевой эйнштейновский квадратичный отклик принят отдельно, первый шаг доказательства уже даёт
\[
G[g](p)\ne0\ \Longrightarrow\ \E[g](p)\ne0,\qquad
\E[g](p)\in\Span\{G[g](p)\}.
\]
Обратное утверждение в нулевой точке получается без полного тождества Уорда. Если $G[g](p)=0$, можно выбрать малые квадратичные возмущения с
$G[g_\epsilon](p)=\epsilon A$, где $A$ - любой заранее заданный ненулевой симметричный тензор. Непрерывность вынуждает
$\E[g](p)\in\Span\{A\}$. Повторив построение с непараллельным $B$, получаем $\E[g](p)=0$.

Следовательно, при соответствующих оставшихся предпосылках
\[
\boxed{\E[g](p)=0\quad\Longleftrightarrow\quad G[g](p)=0.}
\]
Полный Уорд и локальная склейка нужны для более сильного равенства \emph{самих тензоров} $\E=aG$. Это различие существенно: $\E=e^RG$ имеет те же вакуумные нули, но не является сохраняющимся источником для произвольно присоединённой материи. Ни устойчивость динамики, ни постоянство материальной связи из равенства вакуумных нулей не следуют.

\section{Почему плоский квадратичный отклик имеет форму Эйнштейна}
Дифференциал $L=D\E_\eta$ локален, линеен и естественен относительно изометрий $\eta$. На квадратичных нормальных возмущениях это эквивариантное линейное отображение алгебраических тензоров кривизны в симметричные двухтензоры. Возможны только
\[
L_2=A\,\Ric_{\rm lin}+B\,R_{\rm lin}\eta.
\]
Произвольный квадратичный джет отличается от нормального инфинитезимальной заменой координат, задаваемой кубическим векторным полем. При $\E[\eta]=0$ естественность уничтожает эту pure-gauge часть.

Постоянные и линейные вариации не дают кривизны и уничтожаются $L$. Перенос кубического метрического полинома на малый вектор меняет его на квадратичный полином плюс члены степени не выше единицы. Поэтому пространственный дифференциал $L$ на таком кубическом тесте определяется именно $L_2$. Линеаризация Уорда и линейное тождество Бианки дают
\[
(A/2+B)\,\dd R_{\rm lin}=0.
\]
Поскольку кубические вариации допускают произвольный $\dd R_{\rm lin}$, имеем $B=-A/2$. Значит $L_2=aG_{\rm lin}$. Нелинейный оператор не обязан иметь заранее ограниченный порядок производных.

\subsection{Калибровка ненулевости одним существующим свидетелем}
Рассмотрим постоянный квадратичный Brinkmann-профиль
\[
g_A=2\,\dd u\,\dd v+(x^\top Ax)\dd u^2-\dd x_1^2-\dd x_2^2,
\qquad \Tr A\ne0.
\]
Он принадлежит существующему положительному семейству D0 при постоянной $\Psi'(u)$. Предполагается, что этот свидетель и всё используемое семейство бустов принадлежат одной допустимой области источника. Нулевой буст
\[
B_t(u,v,x_1,x_2)=(tu,t^{-1}v,x_1,x_2)
\]
удовлетворяет $B_t^*g_A=\eta+t^2h_2$. По естественности и плоской дифференцируемости
\[
B_t^*\E[g_A](0)=t^2D\E_\eta[h_2](0)+o(t^2).
\]
У ковариантного двухтензора бустовые веса равны $2,1,0,-1,-2$; лишь компонент $uu$ имеет вес $2$. Компоненты меньшего веса должны исчезнуть. Деление компоненты $uu$ на $t^2$ даёт точное $\E[g_A](0)=aG[g_A](0)$.

Если \emph{тот же физический источник} действительно допускает
$\E[g_A]=\kappa T[g_A]\ne0$, то $a\ne0$. Совпадение для ранее сконфигурированного математического источника само по себе не является такой физической калибровкой.

\section{Контрпримеры, которые защищают утверждение}
\subsection{Чистый и смешанный вклад $R^2$}
Пусть
\[
H^{(R^2)}_{\mu\nu}=
2R\Ric_{\mu\nu}-\tfrac12R^2g_{\mu\nu}
+2(g_{\mu\nu}\Box-\nabla_\mu\nabla_\nu)R.
\]
Это локальный естественный сохраняющийся Euler-тензор действия $R^2$. Он удовлетворяет
$H^{(R^2)}[\rho g]=\rho^{-1}H^{(R^2)}[g]$, но его нормальный квадратичный плоский коэффициент равен нулю. Поэтому условие $a\ne0$ необходимо.

Для непосредственной проверки смешанного источника использована метрика
\[
g=e^{2\sigma}\operatorname{diag}(1,-1,-1,-1),\qquad
\sigma=b x^2+c_4x^4.
\]
В начале координат символьный расчёт кривизны даёт
\begin{align*}
G&=\operatorname{diag}(-4b,0,4b,4b),\\
H^{(R^2)}&=\operatorname{diag}(-24b^2-288c_4,-72b^2,
24b^2+288c_4,24b^2+288c_4).
\end{align*}
Для $\E=G+\alpha H^{(R^2)}$ один проективный минор равен
\[
\E_{00}[g]\E_{11}[\rho g]-\E_{11}[g]\E_{00}[\rho g]
=-\frac{288\alpha b^3(\rho-1)}{\rho}.
\]
При $\alpha b\ne0$ и $\rho\ne1$ он ненулевой. Проективный критерий действительно обнаруживает эту поправку, хотя на прежнем Brinkmann-семействе $R^2$ может оставаться невидимым.

\subsection{Без Уорда остаётся переменный множитель}
Источник $\E=e^RG$ локален, гладок, естественен и удовлетворяет проективному закону:
\[
\E[\rho g]=e^{(\rho^{-1}-1)R[g]}\E[g].
\]
Однако
\[
\nabla^\mu\E_{\mu\nu}=e^RG_{\mu\nu}\nabla^\mu R
\]
обычно не равен нулю. Для проверенного конформного джета
$\sigma=b(t^2+x^2)+ct^3$ соответствующий компонент дефекта в начале координат равен $144bc$.

\subsection{Без локальности даже настоящий Euler-источник не жёсток}
На допустимой компактной области с фиксированными граничными данными положим
\[
I[g]=\int R_g\,\dd V_g,\qquad
S[g]=I[g]+\tfrac13 I[g]^3.
\]
С точностью до общей конвенции знака его источник равен
\[
\E[g]=(1+I[g]^2)G[g].
\]
Для каждой фиксированной метрики множитель постоянен в пространстве, поэтому Уорд выполнен. В четырёх измерениях $I[\rho g]=\rho I[g]$, следовательно, закон проективен, а плоский коэффициент равен единице. Но изменение далёкой части метрики меняет источник в данной точке. Нарушены зависимость от ростка и совместимость с ограничением области. Переименование глобального множителя полной истории не заменяет доказательства локального спуска.

\section{Почему малость масштабных ошибок недостаточна}
Для $\rho=\lambda^2$ обозначим
$\Delta[g]=\E[\rho g]-\E[g]$ и $g_n(x)=g(\lambda^{-n}x)$.
Естественность даёт
\[
\lambda^{2(n+1)}\E[g_{n+1}](0)
-\lambda^{2n}\E[g_n](0)
=\lambda^{2n}\Delta[g_n](0).
\]
Телескопирование и плоский предел дают
\[
\boxed{\E[g](p)=aG[g](p)
-\sum_{n=0}^{\infty}\lambda^{2n}\Delta[g_n](0).}
\]
Ряд здесь сначала понимается как предел частичных сумм. Если сумма норм конечна, она ограничивает величину \emph{всей} поправки к Эйнштейну. Она не делает поправку нулевой. Источник $G+\alpha H^{(R^2)}$ имеет геометрически суммируемые дефекты и сохраняет ненулевой $R^2$-вклад. Поэтому этот диагностический критерий нельзя подменять проективной жёсткостью.

\section{Конструктивная попытка по C: квантовый перенос Дуба}
Проверена попытка связать реальный полный процесс со взвешенным терминальным объектом, сохранив когерентную память.

Пусть $V:\mathcal H_0\to\mathcal H_1$ - полная инициализированная изометрия, $E>0$ - терминальный вес и $H=V^\dagger EV$. Тогда
\[
B=E^{1/2}VH^{-1/2},\qquad B^\dagger B=I.
\]
Это точная новая изометрия. Для последовательности полных шагов $W_j$:
\[
H_N=E,\quad H_j=W_j^\dagger H_{j+1}W_j,\quad
B_j=H_{j+1}^{1/2}W_jH_j^{-1/2}.
\]
При положительности требуемых $H_j$ произведение телескопируется:
\[
B_{N-1}\cdots B_0=
E^{1/2}(W_{N-1}\cdots W_0)H_0^{-1/2}.
\]
Все когерентные перекрёстные члены остаются. Абсолютная масса
$Z=\Tr(\rho H_0)$ сохраняет собственный источник
\[
\dd\log Z=\frac{\Tr(\dd\rho\,H_0)+\Tr(\rho\,\dd H_0)}{Z}.
\]

\subsection{Решающий тест на том же возврате}
Для исходного born-return проектора $P$ берём
\[
E=(1-z)I+zP,\qquad 0<z<1.
\]
На скалярной секции $V^\dagger PV=qI$, $0<q<1$, получаем
\[
H=(1-z+zq)I,\qquad
B^\dagger PB=\frac{q}{1-z+zq}I.
\]
Вероятность реальной информативной метки меняется на
\[
\boxed{q_{\rm Doob}-q=
\frac{zq(1-q)}{1-z+zq}>0.}
\]
Поэтому это не перенос цены исходной истории с сохранённым условием задачи. Это другой процесс; его возможность реализовать ещё не доказывает C для прежнего процесса.

Точный общий критерий:
\[
B=V\quad\Longleftrightarrow\quad [E,VV^\dagger]=0.
\]
Если образ $V$ сохраняется положительным $E$, функциональное исчисление даёт
$E^{1/2}V=VH^{1/2}$ и $B=V$. Обратно, $B=V$ означает сохранение этого образа оператором $E^{1/2}$, а значит и $E$; самосопряжённость обеспечивает редуцирующее подпространство.

Семь точных символьных проверок подтверждают конечную алгебру этого теста. Это не общий запрет на D0 и не доказательство невозможности другого физического переноса.

\section{Что теперь требуется для физического применения}
Новый результат позволяет не вычислять заранее неизвестный коэффициент масштабирования \emph{действительного} источника. Но он не предоставляет этот источник автоматически.

Для C по-прежнему нужен один допустимый полный процесс и его реальная цена, с включёнными подготовкой, работой, часами, записями, нормировками и реакциями ограничений. После законного исключения внутренних переменных необходимо доказать свойства того же источника, перечисленные в теореме.

В частности:
\begin{enumerate}
\item Спуск должен относиться ко всем допустимым вариациям, а не только к выбранной волновой семье.
\item Проективное масштабирование - утверждение о полном источнике; независимый дополнительный источник работы может его нарушить.
\item Ненулевой отклик проверяется для той же физической цены, а не переносится по имени из условного $H+J$.
\item В системе с материей сохраняются все её Euler-строки и общий коэффициент связи.
\item Для конечных приближений сохраняются настоящий gain $\Gamma_h$ и норма правого подъёма $R_h$. Если нормированная ошибка равна $\varepsilon_h$, недостаточно одной $\varepsilon_h\to0$: требуется заявленный контроль $\|R_h\|\|\varepsilon_h\|\to0$.
\end{enumerate}

Научный итог этой сессии - доказанное усиление глобального математического критерия и конкретная проверенная попытка по C. Физическое закрытие C и всего D0 $\to$ ОТО в этой сессии не установлено.

\section{Воспроизводимость и состав пакета}
Архив содержит полный английский текст аналитического доказательства, эту записку и её \LaTeX-исходник, два Python-скрипта, наблюдавшиеся JSON-результаты, независимую рецензию и сведения об исходных версиях.

Для повторения конечных проверок нужны Python~3 и SymPy:
\begin{verbatim}
python check_projective_rigidity.py
python doob_actual_mark.py
\end{verbatim}
Каждый скрипт пишет JSON рядом с собой. Проверки вычисляют кривизну, миноры, дефекты и конечные матричные тождества. Они не являются машинным доказательством теоремы о произвольных гладких метриках и не выбирают физическую цену.

\begin{thebibliography}{9}
\bibitem{navarro}
J. Navarro, J. B. Sancho.
\emph{On the naturalness of Einstein's equation}.
Journal of Geometry and Physics 58 (2008), 1007--1014.
\url{https://arxiv.org/abs/0709.1928}.

\bibitem{golden}
D0, \emph{Golden discrete blow-up and condensed GR}, 20.10.2026.
Исследовательский документ о дискретном нормальном сжатии,
каталог \texttt{02\_REGISTRY/research}, commit \texttt{dc411daf}.
\href{https://github.com/gvakhrushev/d0_15/blob/dc411daf01f697ec85dd22f28bfe124171796645/02_REGISTRY/research/A4D_GOLDEN_DISCRETE_BLOWUP_CONDENSED_GR_2026_10_20.md}{Постоянная ссылка на проверенную версию}.

\bibitem{primary}
D0, основной source owner, \S7fu, matched-source synthesis и
\emph{Actual native amplitude representation and terminal-source substitution}.
Основной документ о полном источнике на одном представителе,
commit \texttt{dc411daf}.
\href{https://github.com/gvakhrushev/d0_15/blob/dc411daf01f697ec85dd22f28bfe124171796645/02_REGISTRY/research/A4D_NONVACUUM_SAME_REPRESENTATIVE_PAID_SOURCE_RESIDUAL_2026_10_20.md}{Постоянная ссылка на проверенную версию}.

\bibitem{contract}
D0, \emph{Краткий contract scoped native $\kappa$ $\to$ Einstein closure},
03.11.2026, commit \texttt{dc411daf}.
\href{https://github.com/gvakhrushev/d0_15/blob/dc411daf01f697ec85dd22f28bfe124171796645/02_REGISTRY/research/A4D_SCOPED_NATIVE_KAPPA_CONTRACT_2026_11_03.md}{Контракт A--D}.

\bibitem{book}
D0, Book03, \S\S03.3.0--03.3.0a и 03.25,
commit \texttt{dc411daf}.
\href{https://github.com/gvakhrushev/d0_15/blob/dc411daf01f697ec85dd22f28bfe124171796645/01_BOOKS/BOOK_03_FINITE_ACTION_OPERATORS_AND_SCENE_DYNAMICS.md}{Определения историй, Cost--Info и bootstrap}.
\end{thebibliography}
\end{document}

